\begin{abstract}
This paper sets up a Ramsey growth model with exhaustible natural resources to study the optimal recycling of polluting raw materials and household waste products. During the process of economic development it is socially optimal for the economy to go through an initial “linear” phase with no recycling followed by a “circular” phase where some materials and waste products are recycled to alleviate growing natural resource scarcity and environmental degradation. We show that Pigouvian taxes on polluting activities are not sufficient to ensure the optimal degree of recycling in a market economy; they must be supplemented by subsidies to waste collection and recycling. Numerical simulations of our model illustrate the importance of technical progress and government intervention in the various economic sectors for the transition to a circular economy.

\bigskip\noindent
\textbf{Keywords:} Circular economy, linear economy, optimal recycling, Pigouvian taxation


\bigskip\noindent
\textbf{JEL Classification:}
Q53 {\textbullet}
Q58 {\textbullet}
H21
\end{abstract}